\chapter*{Preface}
This book follows one mathematical statement from its definitions to the laboratory. The statement says that the
triangular lattice has the least energy among all planar configurations of a given density, for every completely
monotone repulsion. The book's line is a single one. A human proof explains \emph{why} a statement holds. A Lean~4 proof
checks that the stated proof is correct, relative to the stated definitions. A numerical experiment can only show that
nothing contradicts the statement in the cases tried. A physical system shows where the mathematics applies and where
it does not.

Every claim carries a tier from the SocrateAI-Scientific-Elenchus method~\cite{elenchus}:
\begin{itemize}
\item \tierA{} kernel-checked Lean;
\item \tierL{} cited literature;
\item \tierC{} argument, conjecture or proposal;
\item \tierX{} floating-point exploration, which may steer a search but never supports a claim.
\end{itemize}
Read the tier before the sentence.

Three things are \emph{not} done here, and the book says so wherever they come up:
\begin{itemize}
\item the strong (``per-potential'') Cohn--Kumar conjecture;
\item a formal treatment of renormalized energies;
\item the Riesz endpoint $s=2$.
\end{itemize}
Renormalized energies are proved in upstream's manuscript, \S8 (\tierL{}), but not formalised anywhere; the strong
per-potential conjecture and the endpoint $s=2$ are open. Chapter~5 explains what each would require.
